\section{Flat-Stable LoRA}
\label{sec:method}

We propose Flat-Stable LoRA (FS-LoRA), a decoupled optimizer for replay-free continual LoRA adaptation.
The method is motivated by a simple design issue: current-task flatness control and past-task Fisher stabilization play different roles, but a naive sharpness-aware optimizer applied to their summed objective mixes these two signals inside the same perturbation closure.
FS-LoRA avoids this coupling by computing the sharpness-aware correction only from the current-task loss, while injecting the Fisher gradient as an independent stability component in the effective adapter-update space $\Delta W=BA$.

The method has two active components.
First, at every mini-batch, FS-LoRA performs a decoupled flat-stable update that combines the clean current-task gradient, the AS$^1$ flatness correction, and the Fisher stability gradient.
Second, at task boundaries, FS-LoRA computes a trace-normalized Fisher in $\Delta W$-space, which prevents the Fisher term from becoming structurally inactive under low-rank parameterization.
An optional dual-ascent controller can adapt $\lambda_{\mathrm{ewc}}$ when the Fisher drift constraint is violated, although it remains passive in the experiments reported here.

\subsection{Why Coupled Sharpness-Aware Fisher Regularization Is Problematic}
\label{sec:coupled_problem}

A natural baseline is to apply GAM or SAM to the combined objective
\begin{equation}
\min_{\phi}
\;
L_t(\phi)
+
\lambda_{\mathrm{ewc}}
L_{\mathrm{Fisher}}(\phi),
\label{eq:coupled_objective}
\end{equation}
where $L_t$ is the current-task loss and $L_{\mathrm{Fisher}}$ penalizes drift along old-task sensitive directions.
Under this design, the perturbation direction used by the sharpness-aware optimizer becomes
\begin{equation}
\hat{\phi}_{\mathrm{coupled}}
=
\phi
+
\rho
\frac{
\nabla_{\phi}
\left(
L_t(\phi)
+
\lambda_{\mathrm{ewc}}
L_{\mathrm{Fisher}}(\phi)
\right)
}{
\left\|
\nabla_{\phi}
\left(
L_t(\phi)
+
\lambda_{\mathrm{ewc}}
L_{\mathrm{Fisher}}(\phi)
\right)
\right\|
+
\epsilon
}.
\label{eq:coupled_perturb}
\end{equation}

This coupled closure is undesirable because the perturbation no longer measures current-task sharpness alone.
To see this, consider a local first-order expansion around $\phi$.
Let $\epsilon_{\mathrm{c}}=\hat{\phi}_{\mathrm{coupled}}-\phi$.
Then
\begin{equation}
\nabla_{\phi} L_t(\phi+\epsilon_{\mathrm{c}})
\approx
\nabla_{\phi} L_t(\phi)
+
H_t \epsilon_{\mathrm{c}},
\label{eq:task_expansion}
\end{equation}
where $H_t=\nabla^2_{\phi}L_t$ is the current-task Hessian.
Similarly, the Fisher penalty contributes
\begin{equation}
\nabla_{\phi} L_{\mathrm{Fisher}}(\phi+\epsilon_{\mathrm{c}})
\approx
\nabla_{\phi} L_{\mathrm{Fisher}}(\phi)
+
H_{\mathrm{Fisher}}\epsilon_{\mathrm{c}}.
\label{eq:fisher_expansion}
\end{equation}
Since the Fisher penalty is defined in the effective update space $\Delta W=BA$, its Hessian in the factor space is locally of the form
\begin{equation}
H_{\mathrm{Fisher}}
\approx
J^\top
\widetilde F
J,
\label{eq:fisher_hessian}
\end{equation}
where $J=\partial \mathrm{vec}(\Delta W)/\partial\phi$ is the Jacobian of the LoRA reparameterization and $\widetilde F$ is the trace-normalized Fisher.

Therefore, applying GAM to the coupled objective injects a Fisher-aligned curvature term
\begin{equation}
J^\top
\widetilde F
J
\epsilon_{\mathrm{c}}
\label{eq:fisher_cross_term}
\end{equation}
into the flatness correction.
The resulting update no longer separates current-task plasticity from old-task stability.
In other words, the sharpness correction is contaminated by the Fisher penalty, and it becomes unclear whether the update improves current-task flatness, protects old-task sensitive directions, or merely mixes the two.

FS-LoRA avoids this failure mode by decoupling the two computations: the sharpness-aware perturbation is computed only from $L_t$, while Fisher regularization is added as a separate stability gradient.

\subsection{Decoupled Flat-Stable Update}
\label{sec:decoupled_update}

Let $\phi$ denote all trainable LoRA factors.
For each mini-batch, FS-LoRA first computes the clean current-task gradient
\begin{equation}
g_{\mathrm{clean}}
=
\nabla_{\phi} L_t(\phi).
\label{eq:gclean}
\end{equation}
The GAM perturbation is then computed using only the current-task loss:
\begin{equation}
\hat{\phi}
=
\phi
+
\rho
\frac{
g_{\mathrm{clean}}
}{
\|g_{\mathrm{clean}}\|+\epsilon
}.
\label{eq:task_only_perturb}
\end{equation}
The perturbed current-task gradient is
\begin{equation}
g_{\mathrm{gam}}
=
\nabla_{\phi}L_t(\hat{\phi}),
\label{eq:ggam}
\end{equation}
and the AS$^1$ flatness correction is defined as the gradient displacement
\begin{equation}
g_{\mathrm{flat}}
=
g_{\mathrm{gam}}
-
g_{\mathrm{clean}}.
\label{eq:gflat}
\end{equation}
Locally, this term satisfies
\begin{equation}
g_{\mathrm{flat}}
\approx
H_t
\left(
\rho
\frac{
g_{\mathrm{clean}}
}{
\|g_{\mathrm{clean}}\|+\epsilon
}
\right),
\label{eq:gflat_hessian}
\end{equation}
so it measures how the current-task gradient changes along the sharpness-aware perturbation direction.
This is the operational AS$^1$ correction used by FS-LoRA.

The Fisher stability gradient is computed independently from the current-task GAM closure:
\begin{equation}
g_{\mathrm{fisher}}
=
\nabla_{\phi}
\left[
\frac{\lambda_{\mathrm{ewc}}}{2}
\mathcal R_{\mathrm{Fisher}}(\phi)
\right].
\label{eq:gfisher}
\end{equation}
The final FS-LoRA update is
\begin{equation}
g_{\mathrm{FS}}
=
g_{\mathrm{clean}}
+
\lambda_{\mathrm{flat}}g_{\mathrm{flat}}
+
g_{\mathrm{fisher}}.
\label{eq:gfs}
\end{equation}

This update is not a generic sum of losses.
It implements a functional separation between plasticity and stability: $g_{\mathrm{flat}}$ is computed from the current-task loss only, while $g_{\mathrm{fisher}}$ constrains the adapter drift along old-task sensitive directions.
The cosine between the two non-clean components,
\begin{equation}
\cos_{\mathrm{flat,fisher}}
=
\frac{
\langle g_{\mathrm{flat}},g_{\mathrm{fisher}}\rangle
}{
\|g_{\mathrm{flat}}\|\|g_{\mathrm{fisher}}\|+\epsilon
},
\label{eq:cos_flat_fisher}
\end{equation}
is logged as a diagnostic.
Empirically, this cosine remains close to zero across the evaluated benchmarks, suggesting that the two channels capture distinct update directions.
We use this near-orthogonality as a design diagnostic rather than as an assumption required by the algorithm.

\subsection{Trace-Normalized Fisher in Effective Adapter-Update Space}
\label{sec:trace_normalized_fisher}

The Fisher penalty is defined in the effective adapter-update space rather than in the raw factor coordinates.
For each LoRA layer,
\begin{equation}
\Delta W_l
=
B_lA_l.
\label{eq:delta_w}
\end{equation}
At the end of task $t'$, FS-LoRA stores the effective adapter snapshot
\begin{equation}
\Delta W_{t',l}^{\star}
=
B_{t',l}^{\star}A_{t',l}^{\star}.
\label{eq:delta_snapshot}
\end{equation}
For a later task $t>t'$, the Fisher drift penalty is
\begin{equation}
\mathcal R_{\mathrm{Fisher}}(\phi)
=
\sum_{t'<t}
\sum_l
\left\langle
\widetilde F_{t',l}+\eta,
\left(
\Delta W_l-\Delta W_{t',l}^{\star}
\right)^2
\right\rangle,
\label{eq:fisher_penalty}
\end{equation}
where $\langle\cdot,\cdot\rangle$ denotes elementwise inner product and the square is elementwise.
This formulation regularizes the functional update $\Delta W=BA$, not the arbitrary factorization $(A,B)$.

The raw Fisher in $\Delta W$-space is often extremely small because the trainable update is restricted to a low-rank subspace.
In practice, raw Fisher values can be around $10^{-6}$, which makes the Fisher gradient nearly inactive unless $\lambda_{\mathrm{ewc}}$ is heavily tuned.
To remove this scale pathology, FS-LoRA uses trace normalization:
\begin{equation}
\widetilde F_{t,l}
=
\frac{
F_{t,l}
}{
\sum_{l'}
\operatorname{tr}
\left(
F_{t,l'}
\right)
+
\epsilon
}.
\label{eq:trace_norm}
\end{equation}
Thus,
\begin{equation}
\sum_l
\operatorname{tr}
\left(
\widetilde F_{t,l}
\right)
\approx
1.
\label{eq:normalized_trace}
\end{equation}
After normalization, $\lambda_{\mathrm{ewc}}$ controls the relative strength of the Fisher constraint rather than compensating for dataset-dependent Fisher scale.

The gradient of Eq.~\eqref{eq:fisher_penalty} follows from the chain rule through $\Delta W_l=B_lA_l$.
Let
\begin{equation}
D_{t',l}
=
\Delta W_l
-
\Delta W_{t',l}^{\star}.
\label{eq:fisher_drift}
\end{equation}
Then
\begin{equation}
\frac{
\partial
\mathcal R_{\mathrm{Fisher}}
}{
\partial B_l
}
=
2
\sum_{t'<t}
\left[
\left(
\widetilde F_{t',l}+\eta
\right)
\odot
D_{t',l}
\right]
A_l^\top,
\label{eq:fisher_grad_b}
\end{equation}
and
\begin{equation}
\frac{
\partial
\mathcal R_{\mathrm{Fisher}}
}{
\partial A_l
}
=
2
\sum_{t'<t}
B_l^\top
\left[
\left(
\widetilde F_{t',l}+\eta
\right)
\odot
D_{t',l}
\right].
\label{eq:fisher_grad_a}
\end{equation}
Because $g_{\mathrm{fisher}}$ in Eq.~\eqref{eq:gfisher} differentiates $\frac{\lambda_{\mathrm{ewc}}}{2}\mathcal R_{\mathrm{Fisher}}$, the factor $2$ in Eqs.~\eqref{eq:fisher_grad_b} and~\eqref{eq:fisher_grad_a} is cancelled by the $\frac{1}{2}$ coefficient.

After each task, the normalized Fisher is accumulated with exponential decay:
\begin{equation}
\bar F_t
\leftarrow
\gamma \bar F_{t-1}
+
(1-\gamma)\widetilde F_t.
\label{eq:fisher_accum}
\end{equation}
This prevents very old tasks from dominating the penalty as the sequence grows.

\subsection{Constrained Optimization View and Optional Dual Ascent}
\label{sec:dual_ascent}

FS-LoRA can be interpreted as a penalty-method approximation to a constrained flat-stable optimization problem.
The desired objective is to improve current-task fit and flatness while keeping the Fisher-weighted drift below a stability budget:
\begin{equation}
\min_{\phi}
\;
L_t(\phi)
+
\lambda_{\mathrm{flat}}
\mathcal S_{\Delta}^{(1)}(\phi)
\quad
\mathrm{s.t.}
\quad
\widetilde D_{\mathrm{Fisher}}(\phi)
\leq
\delta.
\label{eq:constrained_fs}
\end{equation}
Here $\mathcal S_{\Delta}^{(1)}$ denotes the first-order adapter-subspace sharpness surrogate whose gradient is approximated by $g_{\mathrm{flat}}$, and the normalized Fisher drift is
\begin{equation}
\widetilde D_{\mathrm{Fisher}}(\phi)
=
\sum_{t'<t}
\sum_l
\left\langle
\widetilde F_{t',l},
\left(
\Delta W_l-\Delta W_{t',l}^{\star}
\right)^2
\right\rangle.
\label{eq:fisher_drift_budget}
\end{equation}

The corresponding Lagrangian penalty is
\begin{equation}
\mathcal L_{\mathrm{FS}}
(\phi,\lambda_{\mathrm{ewc}})
=
L_t(\phi)
+
\lambda_{\mathrm{flat}}
\mathcal S_{\Delta}^{(1)}(\phi)
+
\frac{\lambda_{\mathrm{ewc}}}{2}
\mathcal R_{\mathrm{Fisher}}(\phi).
\label{eq:fs_lagrangian}
\end{equation}
At each mini-batch, FS-LoRA performs a primal update on $\phi$ using Eq.~\eqref{eq:gfs}.
At task boundaries, an optional dual-ascent controller can update $\lambda_{\mathrm{ewc}}$ according to the observed Fisher drift:
\begin{equation}
\lambda_{\mathrm{ewc}}
\leftarrow
\operatorname{clip}
\left(
\lambda_{\mathrm{ewc}}
+
\alpha_{\mathrm{dual}}
\left(
\widetilde D_{\mathrm{Fisher}}(\phi)
-
\delta
\right),
\lambda_{\min},
\lambda_{\max}
\right).
\label{eq:dual_update}
\end{equation}
When the drift exceeds the budget $\delta$, the Fisher penalty is strengthened for subsequent tasks.
When the constraint is satisfied, $\lambda_{\mathrm{ewc}}$ remains unchanged or decays toward $\lambda_{\min}$, depending on the chosen clipping range.

In the experiments reported in this paper, the Fisher drift remains below the budget, so the dual update is passive.
Therefore, the empirical gains should be attributed to the active components of FS-LoRA, namely the decoupled AS$^1$ update and trace-normalized $\Delta W$-space Fisher stabilization.
Dual ascent is included as an outer-loop controller for regimes where the stability constraint becomes binding, such as task sequences with high semantic overlap.

\subsection{Gradient-Geometric Diagnostic}
\label{sec:gradient_diagnostic}

The decoupled design is motivated by the empirical observation that $g_{\mathrm{flat}}$ and $g_{\mathrm{fisher}}$ are nearly orthogonal in the LoRA factor space.
This observation explains why the two components can be treated as separate update channels, but it is not required by the optimizer.

The connection between the factor space and the effective update space is given by the Jacobian
\begin{equation}
J
=
\frac{
\partial \mathrm{vec}(\Delta W)
}{
\partial \phi
}.
\label{eq:lora_jacobian}
\end{equation}
The Fisher gradient has the form
\begin{equation}
g_{\mathrm{fisher}}
=
J^\top f,
\label{eq:fisher_jacobian_form}
\end{equation}
where $f$ is the Fisher-weighted drift vector in $\Delta W$-space.
The AS$^1$ correction satisfies
\begin{equation}
g_{\mathrm{flat}}
\approx
J^\top
H_{\Delta W}
J
\epsilon,
\label{eq:flat_jacobian_form}
\end{equation}
where $H_{\Delta W}$ is the current-task Hessian in the effective update space and $\epsilon$ is the task-only GAM perturbation.

Thus, the inner product between the two update components is locally
\begin{equation}
g_{\mathrm{flat}}^\top g_{\mathrm{fisher}}
\approx
\epsilon^\top
J^\top
H_{\Delta W}
(JJ^\top)
f.
\label{eq:flat_fisher_inner}
\end{equation}
If the current-task Hessian directions and the old-task Fisher directions are well separated within the reachable update subspace $\operatorname{Im}(J)$, and if $J$ is not highly anisotropic on this subspace, this inner product becomes small.
This provides a geometric explanation for the observed near-zero cosine.

We emphasize that this diagnostic is used to justify the decoupled design, not to claim that orthogonality is enforced by the optimizer.
FS-LoRA does not project gradients, nor does it perform gradient surgery.
It preserves the separation between flatness and stability by avoiding a shared perturbation closure.

\subsection{Algorithm}
\label{sec:algorithm}

\begin{algorithm}[t]
\caption{FS-LoRA Training}
\label{alg:fslora}
\begin{algorithmic}[1]
\Require Task sequence $\mathcal T_1,\ldots,\mathcal T_T$, LoRA parameters $\phi=\{A_l,B_l\}_l$, radius $\rho$, weights $\lambda_{\mathrm{flat}},\lambda_{\mathrm{ewc}}$, decay $\gamma$
\State Initialize LoRA parameters $\phi$ and accumulated Fisher $\bar F_0=0$
\For{$t=1,\ldots,T$}
    \For{mini-batch $\mathcal B\subset \mathcal T_t$}
        \State Compute $g_{\mathrm{clean}}=\nabla_{\phi}L_t(\phi;\mathcal B)$
        \State Compute $\hat{\phi}=\phi+\rho g_{\mathrm{clean}}/(\|g_{\mathrm{clean}}\|+\epsilon)$
        \State Compute $g_{\mathrm{gam}}=\nabla_{\phi}L_t(\hat{\phi};\mathcal B)$
        \State Set $g_{\mathrm{flat}}=g_{\mathrm{gam}}-g_{\mathrm{clean}}$
        \State Compute $g_{\mathrm{fisher}}=\nabla_{\phi}\left[\frac{\lambda_{\mathrm{ewc}}}{2}\mathcal R_{\mathrm{Fisher}}(\phi)\right]$
        \State Update $\phi\leftarrow\phi-\eta\left(g_{\mathrm{clean}}+\lambda_{\mathrm{flat}}g_{\mathrm{flat}}+g_{\mathrm{fisher}}\right)$
    \EndFor
    \State Store $\Delta W_{t,l}^{\star}=B_lA_l$ for all LoRA layers
    \State Estimate $F_{t,l}$ in $\Delta W$-space and compute $\widetilde F_{t,l}$ by trace normalization
    \State Accumulate $\bar F_t\leftarrow\gamma\bar F_{t-1}+(1-\gamma)\widetilde F_t$
    \If{dual ascent is enabled}
        \State Compute $\widetilde D_{\mathrm{Fisher}}(\phi)$
        \State Update $\lambda_{\mathrm{ewc}}$ using Eq.~\eqref{eq:dual_update}
    \EndIf
\EndFor
\end{algorithmic}
\end{algorithm}